\subsubsection{Retorno persistente de las cinco regiones}

Sea
\[
\mathcal F_\pi=\Psi^{-1}(010211)\subset\mathcal U_6
\]
la microfibra relativa al catálogo. Recordemos sus cinco elementos, ahora
separados en bloque inicial y bloque terminal:
\[
\begin{aligned}
 &(501,614,498\mid169,272,272),\\
 &(810,923,870\mid169,272,272),\\
 &(810,983,810\mid169,272,272),\\
 &(870,923,810\mid169,272,272),\\
 &(870,923,810\mid418,674,521).
\end{aligned}
\]
Definimos la proyección de retorno
\[
p_{\rm ret}:\mathbb Z^6\longrightarrow\mathbb Z^3,
\qquad
p_{\rm ret}(u_1,\ldots,u_6)=(u_4,u_5,u_6).
\]

\begin{proposition}[Bloque terminal persistente]
\label{p12-83-c13-s3:prop:bloque-terminal-persistente}
El multiconjunto \(p_{\rm ret}(\mathcal F_\pi)\) posee un único elemento de
multiplicidad máxima:
\[
b_\pi=(169,272,272),
\qquad
\operatorname{mult}_{\mathcal F_\pi}(b_\pi)=4.
\]
El único bloque terminal restante es
\[
b_\pi^-=(418,674,521)
\]
y tiene multiplicidad uno.
\end{proposition}

\begin{proof}
La afirmación se obtiene por enumeración exacta de las cinco filas del
catálogo que satisfacen \(\Psi(U)=010211\). Cuatro filas poseen bloque
terminal \((169,272,272)\); la quinta, que porta el retorno mutado, posee
\((418,674,521)\). La multiplicidad máxima es cuatro y se alcanza una sola
vez.
\end{proof}

La diferencia orientada entre el retorno mutado y el persistente es
\[
 \omega_\pi=b_\pi^- - b_\pi=(249,402,249),
 \qquad
 \langle(1,1,1),\omega_\pi\rangle=900.
\]
Esta identidad separa dos objetos que no deben confundirse. El entero
\(169\) es una coordenada de un dato terminal, es decir, un dato de grado
cero. El vector \(\omega_\pi\) puede leerse como una cocadena de grado uno
después de orientar la arista que une ambos bloques. Llamarlo cociclo exigiría,
además, fijar un complejo de cocadenas y verificar que su coborde se anula.
La prolongación analítica construida más adelante no presupone esa
identificación.

La selección de \(169\) no exige privilegiar la primera posición del bloque.
El grupo simétrico \(\mathfrak S_3\) actúa permutando sus tres coordenadas.
El estabilizador de \(b_\pi\) intercambia las dos apariciones de \(272\) y
fija el único índice de coordenada cuyo valor tiene multiplicidad uno.
Para un bloque de tipo \((r,s,s)\), con \(r\ne s\), denotemos ese índice
por \(\sigma(b)\) y su valor seleccionado por
\[
 \operatorname{sing}(b)=b_{\sigma(b)}.
\]
Para \(\tau\in\mathfrak S_3\), el selector de índices es equivariante
y el valor seleccionado es invariante:
\[
 \sigma(\tau\cdot b)=\tau\bigl(\sigma(b)\bigr),\qquad
 \operatorname{sing}(\tau\cdot b)=\operatorname{sing}(b).
\]
Equivalentemente,
\[
\operatorname{sing}(b)
=\sum_{j=1}^3b_j-2\operatorname{moda}(b).
\]
En particular,
\[
\boxed{\operatorname{sing}(b_\pi)=169.}
\]

\begin{definition}[Selector residual de la microfibra]
\label{p12-83-c13-s3:def:selector-residual-pi}
El selector \(\rho_\pi\) aplica sucesivamente dos operaciones:
\begin{enumerate}
  \item selecciona el bloque terminal de multiplicidad máxima dentro de
  \(p_{\rm ret}(\mathcal F_\pi)\);
  \item selecciona el índice de coordenada fijo por el estabilizador de
  ese bloque y retiene el valor correspondiente.
\end{enumerate}
Por la proposición~\ref{p12-83-c13-s3:prop:bloque-terminal-persistente},
\[
\boxed{\rho_\pi=169.}
\]
\end{definition}

\begin{theorem}[Unicidad equivariante del retorno]
\label{p12-83-c13-s3:thm:unicidad-retorno-169}
Entre los selectores que
\begin{enumerate}
  \item son invariantes bajo permutaciones de las cinco regiones;
  \item seleccionan un índice de coordenada de manera equivariante bajo
  permutaciones de las tres coordenadas terminales;
  \item seleccionan el bloque de multiplicidad máxima;
  \item seleccionan un índice de coordenada fijo por su estabilizador,
\end{enumerate}
el valor de retorno está determinado de manera única y es \(169\).
\end{theorem}

\begin{proof}
La invariancia bajo \(\mathfrak S_5\) impide usar el orden de las filas. La
tercera condición selecciona \(b_\pi\), que es el único modo por la
proposición~\ref{p12-83-c13-s3:prop:bloque-terminal-persistente}. Su estabilizador en
\(\mathfrak S_3\) posee dos órbitas sobre los índices de coordenada:
la pareja que porta el valor \(272\) y el índice unipuntual que porta
\(169\). La cuarta condición selecciona la segunda órbita. La
equivarianza transporta el índice al permutar las coordenadas; el
valor portado por ese índice permanece igual a \(169\).
\end{proof}

Este argumento mejora la lectura posicional del retorno. El \(169\) no es
«la primera coordenada» de una fila escogida: es un invariante del
multiconjunto regional y de la acción de sus simetrías de reordenación.
Tampoco se lo selecciona por rareza global. En las \(468\) emisiones del
catálogo, \(169\) aparece \(84\) veces, uniformemente distribuido con
\(14\) apariciones en cada una de las seis posiciones. La unicidad del
teorema~\ref{p12-83-c13-s3:thm:unicidad-retorno-169} es, por tanto, relacional y regional: reside
en la combinación de microfibra, persistencia modal y estabilizador, no en la
frecuencia aislada del entero.

\subsubsection{Longitud cilíndrica y grado analítico}

Sea \(\mathcal P\) el módulo libre generado por las palabras regionales y
sea \(F^n\mathcal P\) su filtración por longitud. La graduación asociada
registra el primer nivel en el que aparece cada palabra. Para un parámetro
formal \(z\), el carácter ordinario de longitud es
\[
\chi_z([w])=z^{|w|}.
\]
La multiplicatividad
\[
\chi_z([uv])=\chi_z([u])\chi_z([v])
\]
expresa que la concatenación suma longitudes. Una palabra de longitud
\(n\) tiene grado formal \(n\), representado por el monomio \(z^n\).
La evaluación posterior en \(z=\alpha\) asigna a ese monomio el peso
numérico \(\alpha^n\).

\begin{definition}[Compatibilidad de grado]
\label{p12-83-c13-s3:def:compatibilidad-grado}
El lector analítico regional es compatible con la filtración cilíndrica
cuando envía el componente homogéneo de longitud \(n\) al componente
formal homogéneo de grado \(n\):
\[
\operatorname{gr}_n\mathcal P\longrightarrow
\mathbb R z^n.
\]
La evaluación en \(z=\alpha\) es una operación posterior sobre la
expresión formal.
\end{definition}

Para toda palabra \(w\in\mathcal F_\pi\subset\mathcal U_6\), el dato
común de retorno aparece en grado seis:
\[
\rho_\pi\chi_z([w])=169z^6
\xrightarrow{\;z=\alpha\;}169\alpha^6.
\]
Este término registra una sola vez el impulso común de la microfibra.
Este seis es la longitud de la palabra regional. No es la longitud de los
paseos cerrados del apartado~\ref{regcl27:seccion}. Aquellos paseos recorren a
la vez los cuatro anclajes
\(U_\pi,U_e,U_\varphi,U_{\rm APP}\) y poseen longitud mínima ocho. Una
palabra, una arista del grafo de bloques y un paseo cerrado son objetos de
tipos distintos; la graduación utilizada aquí pertenece al primero.

\subsubsection{El carácter determinantal del transporte común}

El cierre dodecafásico proporciona \(\alpha\). La completación de las tres
parejas nonádicas proporciona la escala \(10^3\), de modo que
\[
A=10^3\alpha.
\]
Elegida la coordenatización angular arquimediana, la coordenada común en radianes es
\[
x=\frac{\pi A}{180}=\frac{50\pi}{9}\alpha.
\]
Sea \(R\) una involución real de traza nula en dimensión dos,
\[
R^2=I_2,
\qquad
\operatorname{tr}R=0,
\]
y consideremos el transporte formal bicapa
\[
\mathcal T(x,y)=\exp(-xI_2-yR).
\]
La variable orientada \(y\) todavía no se ha evaluado. Sin embargo, la acción
de \(\mathcal T\) sobre la línea determinante es independiente de ella:
\[
\begin{aligned}
\det\mathcal T(x,y)
&=\exp\!\left(\operatorname{tr}(-xI_2-yR)\right)\\
&=e^{-2x}.
\end{aligned}
\]
Definimos
\[
\boxed{
D_A=e^{-2x}=e^{-100\pi\alpha/9}.}
\]
El subíndice recuerda que se trata del determinante de la contracción común
asociada a \(A\), no de una función de \(C_*^{\rm HMT}\).

\begin{proposition}[Carácter de la línea determinante]
\label{p12-83-c13-s3:prop:caracter-determinante}
La aplicación
\[
\delta_A:(\mathbb N_0,+)\longrightarrow(\mathbb R_{>0},\cdot),
\qquad
\delta_A(k)=D_A^k,
\]
es el único carácter multiplicativo cuyo valor sobre una prolongación
elemental es \(D_A\).
\end{proposition}

\begin{proof}
La multiplicatividad da
\[
\delta_A(k)
=\delta_A(\underbrace{1+\cdots+1}_{k\text{ veces}})
=\delta_A(1)^k=D_A^k.
\]
La misma identidad prueba existencia y unicidad.
\end{proof}

El determinante permanece invariante bajo conjugación de
\(\mathcal T\) y bajo el intercambio de hojas \(R\mapsto-R\). Por tanto, la
cola construida a partir de \(D_A\) pertenece al modo común. La orientación
aparecerá en el signo con el que el carácter regional corrige la fase.

\Needspace{29\baselineskip}
\subsubsection{Reglas del lector analítico regional}

La extensión analítica se fija mediante las reglas siguientes.

\begin{definition}[Lector regional determinantal]
\label{p12-83-c13-s3:def:lector-regional-determinantal}
El lector regional determinantal satisface:
\begin{enumerate}
  \item \emph{compatibilidad de grado}: la longitud cilíndrica es el grado
  analítico formal;
  \item \emph{descomposición de renovación}: el impulso finito de retorno y
  la continuación estricta pertenecen a sumandos aditivos distintos;
  \item \emph{normalización de continuación}: después de registrar una sola
  vez el residuo \(\rho_\pi\), la rama superviviente comienza en la unidad de
  la línea determinante;
  \item \emph{covariancia por concatenación}: cada prolongación elemental
  actúa mediante
  \(\bigwedge^2\mathcal T=D_A\,\operatorname{Id}_{\bigwedge^2\mathbb R^2}\);
  \item \emph{orientación}: en la convención cardinal fijada para APP, el
  retorno regional pertenece al sector compensador y se sustrae de la fase
  de acción de quinto grado.
\end{enumerate}
\end{definition}

La tercera regla tiene contenido matemático. El valor \(169\) se registra
una vez como amplitud del retorno; no se multiplica de nuevo en cada
prolongación. La continuación estricta cuenta la única línea que prosigue y
la normaliza por su unidad. Esta es la diferencia entre una recurrencia de
renovación y la iteración homogénea del impulso inicial.

La necesidad de declarar esa normalización puede probarse. Para todo
\(\lambda\in\mathbb R\), la colección
\[
F_\lambda(z)
=169z^6+\lambda\frac{D_Az^7}{1-D_Az}
\]
conserva el mismo retorno finito y la misma recurrencia multiplicativa de
los coeficientes de la cola. El catálogo y esa recurrencia por sí solos no distinguen
\(\lambda\). La unidad de la línea determinante fija
\(\lambda=1\) antes de evaluar \(z=\alpha\). De este modo, la normalización
no se obtiene ajustando el valor final; forma parte de la definición del
carácter.
